%%%%%%%%%%%%
%
% $Autor: Wings $
% $Datum: 2019-03-05 08:03:15Z $
% $Pfad:  Wings-hub\210520PortentaH7\PortentaH7\$
% $Name: PortentaH7$
% $Version: 4250 $
%
%%%%%%%%%%%%




\chapter{Entwicklungumgebung OpenMV}


\section{Einführung}

%>>>>>>>>>>>>> bearbeiten


Die OpenMV IDE wurde für Machine Vision Anwendungen entwickelt. Sie soll eine Arduino-ähnliche Erfahrung für einfache Computer-Vision-Aufgaben unter Verwendung eines Kamerasensors bieten. OpenMV verfügt über eine eigene Firmware, die auf MicroPython aufbaut. Neben anderer Hardware unterstützt es das Nicla Vision Board. Mit OpenMV können Sie auf einfache Weise eine Vorschau des Kamerastroms anzeigen und Farbbereiche visuell inspizieren, um Schwellenwerte für Ihre Bildverarbeitungsskripte zu definieren. Hier können Sie mehr über die OpenMV IDE lesen.




Die OpenMV IDE ist eine  integrierte Entwicklungsumgebung, die ein Software-Tool zur Verwendung der OpenMV-Kamera eines Arduinos darstellt.  Sie verfügt über einen Text-Editor, einen Frame-Buffer-Viewer mit Histogramm-Anzeige und ein Debug-Terminal. Es kann plattformübergreifend eingesetzt werden und unterstützt Windows, Mac OS, Linux und Raspian. Es wurde für Bildverarbeitungsanwendungen entwickelt und ist in Python programmiert \cite{OpenMV:2021}.  Allerdings setzt das Werkzeug die Verwendung von MicroPython voraus.  \Mynote{Prüfen, ob dies stimmt.} Um MicroPython verwenden zu können, muss die aktuelle Firmware mit dem Bootloader, siehe Kapitel \autoref{chap:Bootloader}


\bigskip

In diesem Abschnitt werden wir die Installation der OpenMV und die erforderlichen Schritte zur Verbindung des Arduino Portenta H7  mit dem Vision Shield beziehungsweise des Arduino Nicla Vision unter Verwendung der OpenMV durchgehen. Außerdem werden wir uns einen Überblick über die Entwicklungsumgebung OpenMV verschaffen und die in der IDE zur Verfügung stehenden Werkzeuge vorstellen.  Schließlich werden wir einige der in der OpenMV bereitgestellten Beispielprogramme implementieren.

%<<<<<<<<<<<<<<<<<< bearbeiten

\section{Installation von OpenMV}

Um die OpenMV IDE zu installieren, muss das Installationsprogramm von der Webseite \textcolor{blue}{\url{https://openmv.io/pages/download}} bzw. \textcolor{blue}{\url{https://github.com/openmv/openmv/releases}} heruntergeladen werden. Der Name der Datei ist \FILE{openmv-ide-windows-2.6.9.exe} und die Größe beträgt 1,21,573 KB. Der Pfad kann nach eigenen Bedürfnissen angegeben werden. 
%Hier habe ich den Pfad als \path{C:\Benutzer\Shoeb\AppData\Roaming\Microsoft\Windows\Startmenü\Programme\OpenMV IDE} eingestellt.

Nachdem der Download abgeschlossen ist, muss die  Setup-Datei geöffnet werden. Es öffnet sich ein Dialog, der darüber informiert, dass die Installation erfolgen kann. Zum Start der Installation muss der Button  \glqq Weiter\grqq{}, siehe \ref{OpenMVInstall}, getätigt werden-

\begin{figure}[H]
	\centering
	\TRANS{
	\includegraphics[width=0.8\textwidth]{Arduino/openmvins1DE}}
    {\includegraphics[width=0.8\textwidth]{Arduino/openmvins1}}
	\caption{Installation von OpenMV}
	\label{OpenMVInstall}
\end{figure}

Nach der Angabe des Installationsordners wird mit dem Betätigen des Buttons \glqq Weiter\grqq{}, wie in der Abbildung \ref{OpenMVInstallFolder} gezeigt, zum nächsten Dialog gewechselt.

\begin{figure}[H]
	\centering
	\TRANS{
	\includegraphics[width=0.8\textwidth]{Arduino/openmvins2DE}}
{\includegraphics[width=0.8\textwidth]{Arduino/openmvins2}}
	\caption{Installationsordner für OpenMV}
	\label{OpenMVInstallFolder}
\end{figure}

Nachfolgend muss der Lizenzvereinbarung, wie in der Abbildung \ref{OpenMVLicence} dargestellt, zugestimmt werden


\begin{figure}[H]
	\centering
	\TRANS{
		\includegraphics[width=0.8\textwidth]{Arduino/openmvins3DE}}
	{\includegraphics[width=0.8\textwidth]{Arduino/openmvins3}}
	\caption{Lizenzverarbeitung für OpenMV}
	\label{OpenMVLicence}
\end{figure}

Jetzt folgt die Wahl für das Kontextmenü für OpenMV, siehe  Abbildung \ref{OpenMVContext}.


\begin{figure}[H]
	\centering
	\TRANS{
		\includegraphics[width=0.8\textwidth]{Arduino/openmvins4DE}}
	{\includegraphics[width=0.8\textwidth]{Arduino/openmvins4}}
	\caption{Kontextmenü für OpenMV}
	\label{OpenMVContext}
\end{figure}


Daraufhin wird ein Dialogfeld mit der Meldung \glqq Bereit zum  Installieren\grqq{} angezeigt. Die Installation startet, wenn wie in Abbildung \ref{OpenMVInstallSetup} \glqq Installieren\grqq{} bestätigt werden.

\begin{figure}[H]
	\centering
	\TRANS{
	\includegraphics[width=0.8\textwidth]{Arduino/openmvins5DE}}
     {\includegraphics[width=0.8\textwidth]{Arduino/openmvins5}}
	\caption{Abschluss der Vorbereitung der Installation von OpenMV}
	\label{OpenMVInstallSetup}
\end{figure}

Während der Installation öffnen sich nacheinander die drei Dialoge \ref{OpenMVDriver}

\begin{figure}[H]
	\centering
	\TRANS{
		\includegraphics[width=0.8\textwidth]{Arduino/openmvins6DE}}
	{\includegraphics[width=0.8\textwidth]{Arduino/openmvins6}}
\quad 
	\TRANS{
	\includegraphics[width=0.8\textwidth]{Arduino/openmvins7DE}}
{\includegraphics[width=0.8\textwidth]{Arduino/openmvins7}}

	\TRANS{
	\includegraphics[width=0.8\textwidth]{Arduino/openmvins8DE}}
{\includegraphics[width=0.8\textwidth]{Arduino/openmvins8}}

	\caption{Installation von Treibern für OpenMV}
	\label{OpenMVDevice}
\end{figure}

Zuletzt öffnet sich der Dialog \ref{OpenMVInstallReady}.

\begin{figure}[H]
	\centering
	\TRANS{
		\includegraphics[width=0.8\textwidth]{Arduino/openmvins9DE}}
	{\includegraphics[width=0.8\textwidth]{Arduino/openmvins9}}
	
	\caption{Fertigstellung der Installation von OpenMV}
	\label{OpenMVInstallReady}
\end{figure}

\begin{figure}[H]
	\centering
	\TRANS{
		\includegraphics[width=0.8\textwidth]{Arduino/openmvins10DE}}
	{\includegraphics[width=0.8\textwidth]{Arduino/openmvins10}}
	
	\caption{Abschluss der Installation und Start von OpenMV}
	\label{OpenMVLaunch}
\end{figure}


Nach der Installation öffnet sich der Dialog \ref{OpenMVLaunch}. Wenn man möchte, wird OpenMV im Anschluss der Installation mit einem Standardprogramm \glqq Hello World\grqq, wie in Abbildung \ref{OpenMVInstallStart} dargestellt, geöffnet.


\begin{figure}[H]
	\centering

	\TRANS{
	\includegraphics[width=0.8\textwidth]{Arduino/openmvideDE}}
{\includegraphics[width=0.8\textwidth]{Arduino/openmvide}}
	\caption{Erster Start von OpenMV}
	\label{OpenMVInstallStart}
\end{figure}


\section{Überblick über OpenMV}

%>>>>>>>>>>>>>>>>>>>>>>>>>

\begin{figure}[H]
	\centering
	
	\begin{tikzpicture}
		
     \node (IDE) at (0,0) {		
		\includegraphics[width=0.8\textwidth]{Arduino/OpenMVIDEStart}};

	 \draw[red] (2.4,1.75) -- (2.4,2.85) -- (5.2,2.85) -- (5.2,1.75) -- cycle; 	
	 \draw[red] (3.6,2.85) -- (3.6,3.5);
	 \node (VS) at (3.0,3.6) {\tiny Videostream};

	 \draw[red] (3.7,1.73) -- (3.7,1.58) -- (5.2,1.58) -- (5.2,1.73) -- cycle; 	
     \draw[red] (5.2,1.65) -- (5.3,1.65);
     \node (VS) at (5.8,1.65) {\tiny Farbraum};


	 \draw[red] (4.6,3.05) -- (4.6,2.85) -- (5.2,2.85) -- (5.2,3.05) -- cycle; 	
	 \draw[red] (4.9,3.05) -- (4.9,4.1);
     \node (AD) at (3.8,4.2) {\tiny Aktivieren/Deaktivieren};
	 
	 \draw[red] (4.22,3.05) -- (4.22,2.85) -- (4.54,2.85) -- (4.54,3.05) -- cycle; 	
     \draw[red] (4.38,3.05) -- (4.38,3.9);
     \node (Z) at (4.2,4.0) {\tiny Zoomen};

	 \draw[red] (3.72,3.05) -- (3.72,2.85) -- (4.18,2.85) -- (4.18,3.05) -- cycle; 	
     \draw[red] (3.95,3.05) -- (3.95,3.7);
     \node (Z) at (3.5,3.8) {\tiny Aufnehmen};

	 \draw[red] (2.4,1.55) -- (2.4,-3.34) -- (5.2,-3.34) -- (5.2,1.55) -- cycle; 	
     \draw[red] (3.8,-3.34) -- (3.8,-3.8);
     \node (VS) at (3.4,-3.9) {\tiny Histrogramm};
	 
	 \draw[darkcyan] (0,-3.53) -- (0,-3.35) -- (5.2,-3.35) -- (5.2,-3.53) -- cycle; 	
     \draw[darkcyan] (1.6,-3.53) -- (1.6,-4);
     \node (VS) at (1.6,-4.1) {\tiny Statuszeile};

	 \draw[darkcyan] (-5.2,-2.85) -- (-5.2,-3.19) -- (-4.7,-3.19) -- (-4.7,-2.85) -- cycle; 	
     \draw[darkcyan] (-5.2,-3) -- (-5.5,-3)  -- (-5.5,-4.1);
     \node (VS) at (-5.0,-4.2) {\tiny Verbinden};

	 \draw[green] (-5.2,-3.52) -- (-5.2,-3.2) -- (-4.7,-3.2) -- (-4.7,-3.52) -- cycle; 	
     \draw[green] (-5.2,-3.4) -- (-5.4,-3.4)  -- (-5.4,-3.8);
     \node (VS) at (-4.3,-3.9) {\tiny Start des Programms};

	 \draw[yellow] (-4.65,-3.54) -- (-4.65,-3.35) -- (-3.0,-3.35) -- (-3.0,-3.54) -- cycle; 	
     \draw[yellow] (-3.1,-3.54) -- (-3.1,-4);
     \node (VS) at (-1.7,-4.1) {\tiny Suchen + Serieller Monitor};


	 \draw[green] (-5.2,1.75) -- (-5.2,2.85) -- (-4.7,2.85) -- (-4.7,1.75) -- cycle; 	
     \draw[green] (-5.2,2) -- (-5.3,2) -- (-5.3,3.6);
     \node (VS) at (-5.05,3.7) {\tiny Datei};

	 \draw[green] (-5.2,0) -- (-5.2,1.73) -- (-4.7,1.73) -- (-4.7,0) -- cycle; 	
     \draw[green] (-5.2,1) -- (-5.4,1) -- (-5.4,4.0);
    \node (VS) at (-4.8,4.1) {\tiny Bearbeitung};


     \draw[green] (-5.2,3.05) -- (-5.2,3.25) -- (-2.7,3.25) -- (-2.7,3.05) -- cycle; 	
     \draw[green] (-3.5,3.25) -- (-3.5,3.6);
     \node (VS) at (-3.0,3.7) {\tiny Menüleiste};
     
	 \draw[pink,line width=0.8pt] (2.3,-3.33) -- (2.3,2.8) -- (-4.6,2.8) -- (-4.6,-3.33) -- cycle; 	
     \draw[pink,line width=0.8pt] (-4.6,1) -- (0,1);
     \node (VS) at (-1,1.2) {\tiny Programmbereich};
     



    \end{tikzpicture}
	\caption[Elemente der OpenMV IDE]{Elemente der OpenMV IDE. Rot: Arbeiten mit dem Videostream. Grün: Arbeiten mit der Datei und die Menüleiste. GrünBlau: Verbinden mit dem Board und Informationen über das Board. Gelb: Zusatzfunktionen. Pink: Der Programmierbereich für MicroPython.}
	\label{OpenMVElements}
\end{figure}

Die Abbildung \ref{OpenMVElements} zeigt die Entwicklungsumgebung, wobei verschiedene Elemente farblich markiert sind. Im Folgenden werden verschiedene Funktionen kurz vorgestellt.

\subsection{Arbeiten mit dem Videostream}


In der Abbildung \ref{OpenMVElements} sind einige Bereiche rot markiert. Diese Bereiche beziehen sich auf den Videostream, der dort dargestellt werden kann. Für den Videostream stehen verschiedene Funktionen zur Verfügung.

\begin{itemize}
  \item  \textbf{Aufnehmen:} Der Videostream kann mit einer Frame von 30 Bilder pro Sekunde aufgezeichnet werden. 
  
   Es kann auch ein einzelnes Bild gespeichert werden, indem einfach mit einem  Klick mit der  rechten Maustaste das Kontextmenü geöffnet wird. Im Kontextmenü muss dann die Option \glqq Sichern als Bild\grqq{} ausgewählt werden.
  \item \textbf{Zoomen:} Mit dieser Option können die Bilder    vergrößert oder verkleinert werden.
  \item \textbf{Deaktivieren:} Mit dieser Option der Videostream zu- oder abgeschaltet werden. Wenn er abgeschaltet ist, wird kein Bild angezeigt.
	
\end{itemize}


OpenMV verfügt über eine Histogramm-Anzeige, die dazu dient, eine Rückmeldung über die Beleuchtungsqualität im Raum zu erhalten, die im Grunde eine Vorstellung von der Qualität des Bildes vermittelt, das die Kamera betrachtet. Zusätzlich stehen hier vier Farbmodelle zur Verfügung:

\begin{itemize}
  \item RGB
  \item YUV
  \item LAB
  \item Graustufen
\end{itemize}  
  
Da in den Programmen die verschiedene Modelle verwendet werden, können mit Hilfe des Histogramms optimal Anpassungen durchgeführt werden.  


\subsection{Zusatzfunktionen}

Die beiden Zusatzfunktionen \glqq Suchen\grqq{} und \glqq Serieller Monitor\grqq{} sind  gelb markiert. Bei der Suchfunktion ist eine Standardfunktion eines Editors, daher wird diese hier nicht erläutert. 



\subsubsection{Serieller Monitor}



Dieses befindet sich am unteren Rand der OpenMV. Der gesamte Debug-Text, der mit dem Befehl \PYTHON{print()} ausgegeben wird, wird hier angezeigt.


%>>>>>>>>>>>>>>>>>>>>> hier weiter:

\subsection{Verbinden mit dem Board und Informationen über das Board}




\textbf{Statusleiste:}

In der Statusleiste zeigt OpenMV die Firmware-Version der OpenMV Cam, den seriellen Anschluss, das Laufwerk und die FPS an. Wir können hier auch die Firmware-Version aktualisieren, wenn sie nicht aktualisiert ist. In der Statusleiste werden die aktuellen FPS, das Board und die Bezeichnung des seriellen Anschlusses angezeigt.


\subsection{Bearbeiten der Datei}




\subsection{Menüleiste}

%todo Hier fehlt noch etwas:

In der Menüleiste befinden sich folgende Punkte:

\begin{itemize}
	\item File
	\item Bearbeiten
	\item Werkzeuge
	\item Fenster
	\item Helfen
\end{itemize}

\subsubsection{Werkzeuge}

Bei den Werkzeugen gibt es viele Optionen. Einge werden im Folgenden erläutert:

\begin{itemize}
	\item Speichern des geöffneten Skripts:
	
	Damit kann das Skript für die OpenMV-Kamera, zum Beispiel dem Arduino-Vision-Shield, gespeichert werden. Dabei werden Kommentare und Leerzeichen automatisch gelöscht, um Speicherplatz zu sparen, und das Skript wird als Python-Datei \FILE{main.py} gespeichert.
	
	\item Konfigurieren der Einstellungsdatei der OpenMV-Kamera:
	
	Dies ermöglicht es, eine \FILE{.ini}-Datei zu ändern, die in der OpenMV-Kamera unter Verwendung von OpenMV gespeichert ist und die die OpenMV-Kamera beim Hochfahren für bestimmte Hardwarekonfigurationen liest.
	
	\item Zurücksetzen der OpenMV-Kamera :
	
	Dieser Befehl setzt die OpenMV-Kamera zurück und trennt die Verbindung.
	
\end{itemize} 


\subsection{Programmbereich  - MicroPython}


\subsubsection{Einführung}

Der Programmbereich steht zur Verfügung,  um Anwendungen in der Programmiersprache MicroPython\index{MicroPython} zu programmieren. Im folgenden Abschnitt werden einige der Funktionen der Programmiersprache MicroPython beschrieben und gezeigt, wie MicroPython-Skripte auf einem Arduino-Board mit Hilfe der OpenMV-IDE ausgeführt werden.

Bevor mit MicroPython begonnen werden kann, muss der Bootloader des Boards\index{Bootloader} aktualisiert werden. Dies geschieht, indem in der Arduino IDE zuerst die neueste Version des \glqq Mbed OS Nicla core\grqq{}  heruntergeladen wird. In der Arduino IDE wird der Sketch \FILE{File $\rightarrow$ Examples $\rightarrow$ STM32H747\_System $\rightarrow$ STM32H747\_updateBootloaderDatei} geöffnet und auf das Board geladen. Nachdem der Sketch hochgeladen wurde, muss den Anweisungen im seriellen Monitor gefolgt werden. Dies ist ausführlich im Abschnitt \ref{chap:Bootloader} beschrieben.

\subsubsection{MicroPython}

MicroPython ist ein Interpreter der Programmiersprache Python, der auf kleinen eingebetteten Entwicklungsplatinen funktioniert und Open Source ist. Anstatt komplexe Low-Level-Sprachen wie C oder C++ zu verwenden, kann mit MicroPython sauberer und einfacher Python-Code geschrieben werden, um die Hardware zu steuern. Es handelt sich um eine schlanke und effiziente Python-3-Implementierung, die nur einen winzigen Teil der Python-Standardbibliothek enthält und für den Einsatz auf Mikrocontrollern und in eingeschränkten Umgebungen optimiert ist. Sie ist kompakt genug, um in nur 256k Code und 16k RAM zu passen und zu laufen.

Es hat einzigartige Eigenschaften, die es für eingebetteten Systeme prädestiniert:

\begin{itemize}
	\item \textbf{Interaktiver REPL (Read-Evaluate-Print-Loop):} Dies hilft bei der Verbindung mit einem Arduino Board und der Ausführung des Codes, ohne dass eine Kompilierung oder ein Upload erforderlich ist. Es handelt sich um einen vierstufigen Prozess, bei dem zunächst die Benutzereingaben gelesen werden, der Code ausgewertet wird, alle Ergebnisse gedruckt werden und eine Schleife zurück zum ersten Schritt erfolgt.
	\item \textbf{Umfangreiche Software-Bibliothek:} MicroPython hat wie die normale Python-Programmiersprache Bibliotheken eingebaut, die viele Aufgaben unterstützen. Das Parsen von JSON-Daten aus einem Webservice, das Durchsuchen von Text mit einem regulären Ausdruck oder sogar das Programmieren von Netzwerk-Sockets ist mit den eingebauten Bibliotheken für MicroPython einfach.
	\item \textbf{Erweiterbarkeit:} MicroPython ist mit C/C++-Funktionen auf niedriger Ebene erweiterbar, d. h. wir können bei Bedarf MicroPython-Code auf hoher Ebene mit schnellerem Code auf niedriger Ebene kombinieren.
\end{itemize}

Die Programmiersprache MicroPython implementiert die Kernsprache Python 3, kann aber nicht die gesamte Standardbibliothek von Python 3 implementieren, da der Versuch, große Bibliotheken in winzige Boards mit nur einem Kilobyte Speicher zu integrieren, nicht möglich ist. Zu beachten ist, dass MicroPython nur ein Interpreter für die Programmiersprache ist, während Arduino ein ganzes Ökosystem ist, das die Arduino IDE und die Hardware, zum Beispiel Arduino Portenta H7 oder Arduino Nicla Vision, umfasst.

MicroPython-Skripte können mit Hilfe der OpenMV-IDE auf dem Arduino Portenta H7 ausgeführt werden \cite{Romero:2020a}. OpenMV erfordert eine eigene Firmware, die auf MicroPython aufbaut. Wir können die Funktionen eines Arduino-Boards testen, indem wir verschiedene in MicroPython bereitgestellte Klassen und Module verwenden.

%Um MicroPython auf einem Arduino-Board zu implementieren, muss zunächst das Arduino-Board über ein USB-C-Kabel mit dem Computer verbinden und ihn durch doppeltes Drücken der Reset-Taste flashen. Nach dem Flashen wird der Bootloader\index{Bootloader} aktualisiert, während die grüne LED blinkt. Nun öffnen wir die OpenMV-IDE und klicken auf die Option \glqq Connect\grqq{} am unteren Rand des Bildschirms, wie in den vorherigen Kapiteln beschrieben. Dadurch wird die neueste OpenMV-Firmware auf dem Arduino-Board installiert. Nachdem die Firmware aktualisiert wurde, müssen wir auf den Button \glqq Play\grqq{} klicken, um das Board mit der OpenMV IDE zu verbinden und mit der Programmierung in der IDE zu beginnen.
%Um zu prüfen, ob ein MicroPython-Programm auf dem Arduino-Board laufen kann, schreiben wir ein einfaches Programm in MicroPython, um die LEDs auf dem Board mit der OpenMV IDE blinken zu lassen.


\subsubsection{Empfohlene Bibliotheken}


\begin{itemize}
	\item  
	\href{https://www.arduino.cc/reference/en/libraries/arduinoble/}{ArduinoBLE}
	
	Die ArduinoBLE-Bibliothek ist für Arduino-Boards konzipiert, deren Hardware für BLE und Bluetooth 4.0 und höher aktiviert ist.
	
	\item \href{https://docs.arduino.cc/learn/built-in-libraries/pdm}{PDM}
	
	Die PDM-Bibliothek ermöglicht die Verwendung von PDM-Mikrofonen (Pulse-Density-Modulation), wie dem MP34DT06JTR.
	
	\item \href{https://www.arduino.cc/reference/en/libraries/wifi/}{Wi-Fi}
	
	Die WiFi-Bibliothek ist für die Verwendung des Murata 1DX-Moduls konzipiert, mit dem Ihr Arduino eine Verbindung zum Internet herstellen kann.
\end{itemize}


\subsubsection{Anleitungen}

Die Firma Arduino stellt einige Anleitungen zur Verfügung, die im Folgenden aufgelistet sind.

\begin{itemize}
	\item \href{https://docs.arduino.cc/tutorials/nicla-vision/blob-detection}{Blob detection with OpenMV}
	
	Dieses Tutorial zeigt Ihnen, wie Sie mit Nicla Vision das Vorhandensein und die Position von Objekten in einem Kamerabild erkennen können.
	
	\medskip
	
	Schlüsselwörter: OpenMV, Blob detection, Machine Vision
	
	\medskip
	
	\item \href{https://docs.arduino.cc/tutorials/nicla-vision/getting-started}{Getting Started with Nicla Vision}
	
	In diesem Tutorial lernen Sie, wie Sie das Board einrichten, die OpenMV IDE verwenden und einen MicroPython-Sketch ausführen.
	
	\medskip
	
	
	Schlüsselwörter: Getting Started, OpenMV, Setup, MicroPython
	
	\medskip
	
	\item \href{https://docs.arduino.cc/tutorials/nicla-vision/image-classification}{Image Classification with Edge Impulse}
	
	In diesem Tutorial lernen Sie, wie Sie mit Edge Impulse ein benutzerdefiniertes maschinelles Lernmodell trainieren und eine Bildklassifizierung mit dem Arduino Nicla Vision durchführen.
	
	\medskip
	
	Schlüsselwörter: Machine Learning, Edge Impulse, TinyML, Tensorflow
	
	\medskip
	
	\item \href{https://docs.arduino.cc/tutorials/nicla-vision/microphone_sensor}{Reading Audio Samples With the Onboard Microphone}
	
	Erfahren Sie, wie Sie mit dem eingebauten Mikrofon des Nicla Vision ein Schallmessgerät erstellen können.
	
	\medskip
	
	Schlüsselwörter:         OpenMV,   Microphone,  Sound,      Sensor
	
	\medskip	
	
	
	\item \href{https://docs.arduino.cc/tutorials/nicla-vision/nicla-vision-imu}{Accessing IMU Data on Nicla Vision}
	
	Erfahren Sie, wie Sie auf die Daten des Beschleunigungsmessers und des Gyroskops zugreifen können, die im Lieferumfang des LSM6DSOXTR IMU-Moduls enthalten sind.
	
	\medskip
	
	Schlüsselwörter:  Gyroscope, Accelerometer
	
	\medskip
	
	\item \href{https://docs.arduino.cc/tutorials/nicla-vision/proximity}{Proximity Detection with Arduino Nicla Vision}
	
	Erfahren Sie, wie Sie mit dem Näherungssensor die Geschwindigkeit des Blinkens der LED variieren können.
	
	\medskip	
	
	Schlüsselwörter: Proximity, Time Of Flight, Blink  
	
\end{itemize}






%<<<<<<<<<<<<<<<<<<<<<<<<<

\section{Verwendung eines Arduino-Boards in OpenMV}

Nachfolgend wird beschrieben, wie das Flashen der OpenMV-Firmware\index{Firmware}\index{Bootloader} durchgeführt werden kann. Dazu wird zuerst das Arduino-Board, zum Beispiel ein Nicla Vision, über ein USB-Kabel an Ihren Computer angeschlossen. Im nächsten Schritt muss das Board in den Bootloader-Modus versetzt werden. Bei Nicla Vision geschieht dies, indem der Reset-Knopf auf der Platine zweimal gedrückt. Die eingebaute LED beginnt zu leuchten. 
\Mynote{Muss Open MV geschlossen sein, oder kann es geöffnet sein?}
%Öffnen Sie nun die OpenMV IDE, siehe \ref{OpenMVStart}.

%\begin{figure}
%  \centering
%  \includegraphics[width=0.8\textwidth]{OpenMV/OpenMVOpenIDE}
%	
%  \caption{OpenMV IDE nach dem Start}\label{OpenMVStart}
%\end{figure}


%\begin{figure}
%	\centering
%	\includegraphics[width=0.8\textwidth]{OpenMV/OpenMVOpenIDE}
%	
%	\caption{OpenMV IDE nach dem Start}
%	\label{OpenMVStart}
%\end{figure}



Dann muss in der OpenMV IDE der Schalter \glqq Verbinden\grqq{}  am unteren Rand der linken Symbolleiste betätigt werden, siehe Abbildung \ref{OpenMVConnect}.
%, in dem  
% \glqq Install the latest release firmware (v4.0.1)\grqq{} angewählt werden soll. Hier wird dies mit  \glqq OK\grqq{} bestätigt, wie in Abbildung \ref{OpenMVInstallFirmware} dargestellt. 
Ein Pop-up-Fenster, siehe Abbildung \ref{OpenMVFirmware}, fragt nach Bestätigung \glqq DFU bootloader(s) found. What would you like to do?\grqq{} zur Installation der neuesten Firmware. Es muss die Option \glqq Reset Firmware to Release Version\grqq{} gewählt werden. Dadurch wird die neueste OpenMV-Firmware auf dem Board installiert. In dem Dialog wird auch gefragt, ob das interne Dateisystem gelöscht werden soll; dies kann verneint werden, siehe \ref{OpenMVFirmware}.\index{Firmware}\index{Bootloader}

\begin{figure}
	\centering
	\includegraphics[width=0.8\textwidth]{OpenMV/OpenMVConnect}
	
	\caption[Verbinden eines Arduinos mit der OpenMV IDE]{Betätigung der Schaltfläche \glqq connect\grqq{}, um den Arduino Portenta oder Nicla Vision mit der OpenMV IDE zu verbinden}
	\label{OpenMVConnect}
\end{figure}


\begin{figure}
	\centering
	\includegraphics[width=0.8\textwidth]{OpenMV/OpenMVFirmware}
	
	\caption{Installation der neuesten Version der OpenMV-Firmware}
	\label{OpenMVFirmware}
\end{figure}



\begin{figure}[H]
	\centering
	\includegraphics[width=0.8\textwidth]{Arduino/openmvinst}
	\caption{Installation der neuesten Firmware}
	\label{OpenMVInstallFirmware}
\end{figure}



Die LED des Boards beginnt zu blinken, während die OpenMV-Firmware hochgeladen wird. Es öffnet sich ein Pop-up-Fenster, das den Fortschritt des Uploads anzeigt. Nach Abschluss der Installation der Firmware hört LED auf zu blinken und verblasst. Wenn der Vorgang abgeschlossen ist, erscheint die Meldung \glqq DFU firmware update complete!\grqq, siehe Abbildung \ref{OpenMVInstallUpdate}. 

\begin{figure}[H]
	\centering
	\includegraphics[width=0.8\textwidth]{Arduino/firmwareinstalled}
	\caption{Abschluss der Firmware-Update}
	\label{OpenMVInstallUpdate}
\end{figure}

\bigskip


Das Nicla Vision beginnt mit seiner blauen LED zu blinken, wenn es bereit ist, verbunden zu werden. Nach Bestätigung des Abschlussdialogs sollte das Board bereits mit der OpenMV IDE verbunden sein, andernfalls muss erneut die Schaltfläche \glqq connect\grqq{},  siehe \ref{OpenMVReady}, betätigt.
Nach Abschluss der Installation kann das Board verwendet werden. Dazu muss ein Programm in der OpenMV IDE geladen sein. Das Programm kann gestartet werden, in dem die  Taste \glqq Play\grqq{}. Wenn das Programm nicht ausgeführt werden kann, müssen wir das Arduino-Board zurücksetzen. Nach dem Zurücksetzen des Boards muss wiederum auf die Schaltfläche \glqq Play\grqq{} betätigt werden. 

\begin{figure}
	\centering
	\includegraphics[width=0.8\textwidth]{OpenMV/OpenMVReady}
	
	\caption[Erfolgreiche Verbindung]{Wenn das Portenta bzw. Nicla Vision erfolgreich mit der OpenMV IDE verbunden ist, erscheint unten links eine grüne Abspieltaste}
	\label{OpenMVReady}
\end{figure}


%%%%%%%%%%%%%%%%%%%%%%%%%%%%%%%%%%%%%


%\textbf{Achtung:} Die Installation der OpenMV-Firmware überschreibt alle vorhandenen Sketches im internen Flash. Infolgedessen wird der Port des Boards in der Arduino IDE nicht mehr sichtbar sein. Um eine Arduino-Firmware neu zu flashen, müssen Sie das Board in den Bootloader-Modus versetzen. Drücken Sie dazu zweimal den Reset-Knopf auf dem Board. Die eingebaute LED fängt an, auf- und abzublenden. Im Bootloader-Modus sehen Sie den Port des Boards wieder in der Arduino-IDE.\index{Firmware}\index{Bootloader}
%
%\bigskip



\section{Beispiele für OpenMV}

\subsection{Hello World \glqq Blinken der eingebauten LED mit OpenMV\grqq}

Im Programm wird zunächst das Modul \PYTHON{pyb} importiert, das Funktionen für die Boards enthält, wie z.B. die PIN-Behandlung. Dann werden Variablen zur Steuerung der eingebauten LEDs mit der Funktion \PYTHON{pyb.LED()} angelegt. In einer \PYTHON{while}-Schleife werden die roten, blauen und grünen LEDs mit einer Verzögerung von 1000ms mit der Funktion \PYTHON{redLED.on()} ein- und schalten die LED mit \PYTHON{redLED.off()} ausgeschaltet. Um das fertige Programm hochzuladen,  muss der Button \glqq Play\grqq{}  unten, um das Skript hochzuladen.


\begin{code}[H]
	\lstinputlisting[language=Python, firstnumber=1]{../Code/MicroPython/Blink.py}
	\caption{Sketch \FILE{Blink.py}}\label{MicroPythonBlink}
\end{code}

Aus der Abbildung \ref{MicroPythonBlinkOutput} geht hervor, dass die Ausgabe in einem Zeitintervall von 1000ms rot, blau und grün blinkt. Damit ist das Programm verifiziert.
Daher können wir sagen, dass MicroPython-Skripte auf dem Arduino-Board unter Verwendung der OpenMV IDE ausgeführt werden können.

\begin{figure}[H]
	\centering
	\includegraphics[width=\textwidth]{Arduino/outputmerged}
	\caption{Ausgang RGB-LED blinkt rot, blau und grün}
	\label{MicroPythonBlinkOutput}
\end{figure}


\subsection{Erstellen eines einfachen Gesichtsfilters mit OpenMV}

OpenMV unterstützt die Bildformate \FILE{pbm}, \FILE{bmp}, \FILE{pgm} oder \FILE{ppm} und nicht die anderen Bildformate wie \FILE{png} \cite{Romero:2020}. In diesem Beispiel wird ein Gesicht im Videostream detektiert und mit einem anderen Bild, in diesem Beispiel ein Smiley-Bild, überlagert. 

Dazu wird zunächst das Bild heruntergeladen und auf einen USB-Stick kopiert.
Im Python-Programm wird zunächst das Bild mit der Funktion \PYTHON{Image()} in eine Variable mit dem Namen \PYTHON{faceImage} geladen. Die Variable \PYTHON{copy\_to\_fb} wird auf den Wert \PYTHON{false} gesetzt; damit wird verhindert dass das Bild automatisch in den Framebuffer kopiert wird. Als Nächstes wird das Skalierungsverhältnis berechnet, um das Bitmap-Bild so zu skalieren, dass es dem erkannten Gesicht im Videostream entspricht, indem die Funktion \PYTHON{scale\_ratio=faceWidth/faceImage.width()} verwendet wird.
Dann zeichnen wir das Bitmap-Bild auf das erkannte Gesicht mit der Funktion \PYTHON{draw\_image()}.

Hier wird die Gesichtserkennung mit dem Haar-Kaskaden-Algorithmus durchgeführt. \Mynote{Quelle?} In diesem Algorithmus gibt es mehrere Stufen, in denen die Ausgabe einer Stufe zusätzliche Informationen für die Eingabe der nächsten Stufe in der Kaskade liefert. Auf diese Weise werden die Kanten, Linien und Kontrastprüfungen in verschiedenen Stufen berechnet, und schließlich wird das Gesicht erkannt und ein Begrenzungsrahmen um es herum gezeichnet. 


\begin{figure}[H]
	\centering
	\includegraphics[width=\textwidth]{Arduino/facefilter}
	\caption{Ausgabe des Gesichtsfilters}
	\label{OpenMVFace}
\end{figure}

Wie aus der Abbildung \ref{OpenMVFace} hervorgeht, wird das zuvor geladene Bitmap-Bild mit dem erkannten Gesicht überlagert, sobald ein Gesicht im Videostream erkannt wird. Das vollständige Programm ist in der Programmumgebung \ref{MicroPythonFaceFilter}.


\begin{code}[H]
	{\tiny 
	\lstinputlisting[language=Python, firstnumber=1]{../Code/OpenMV/FaceFilter.py}
}
	\caption{Erstellen eines einfachen Gesichtsfilters mit OpenMV \FILE{FaceFilter.py}} 	\label{MicroPythonFaceFilter}
\end{code}  


%%%%%%%%%%%%%%%%%%%%%%%%%%%%%%%%%%%%%%%%%%%%%

\subsection{Gesichtsverfolgung}

In der Entwicklungsumgebung OpenMV steht das Beispiel \FILE{face\_tracking}. Mit Hilfe des Programms können detektierte Gesichter nachverfolgt werden. Dazu werden sie mit einem umgebenden Rechteck markiert. 

In dem Programm wird der Sensor zunächst mit der Funktion \PYTHON{sensor.reset()} zurückgesetzt und die Bildgröße auf QVGA und das Pixelformat auf Graustufen eingestellt. Danach wird die Haar-Kaskade mit der Funktion \PYTHON{image.HaarCascade()} geladen. Dann wird die Funktion \glqq Keypoints\grqq{} der OpenMV-Kamera verwendet, um ein Gesicht zu verfolgen, nachdem es durch den Haar-Kaskaden-Algorithmus erkannt wurde. Als erstes wird die Keypoints auf \PYTHON{none} gesetzt und, wenn das Gesicht erkannt wird, wird die Funktion \PYTHON{image.find\_keypoints()} verwendet. Wenn das Gesicht erkannt wurde, wird mit der Funktion \PYTHON{draw\_rectangle()} ein Rechteck um das erkannte Gesicht  und die Keypoints mit der Funktion \PYTHON{draw\_keypoints()} gezeichnet. Danach werden die Keypoints im gesamten Frame extrahiert und diese Keypoints mit den ersten Keypoints abgeglichen, die nur für das Gesicht mit der Funktion \PYTHON{image.match\_descriptor()} verwendet wurden. Wenn die Übereinstimmung der Keypoints größer als 5 ist, wird ein Rechteck um das übereinstimmende Merkmal, siehe \ref{OpenMVFaceTracking}, gezeichnet.

\begin{figure}[H]
	\centering
	\includegraphics[width=\textwidth]{Arduino/face_tracking}
	\caption{Ausgabe der Gesichtsverfolgung}
	\label{OpenMVFaceTracking}
\end{figure}

Zum Test des Sketchs muss er ausgeführt werden. Während der Ausführung muss die Kamera auf ein Gesicht gerichtet und richtig fokussiert werden. Sobald das Gesicht im Videostream erkannt wird, wird ein rechteckiger Rahmen um das Gesicht herum gebildet, der das erkannte Gesicht im Bild anzeigt. 

\begin{code}[H]
	\lstinputlisting[language=Python, firstnumber=1]{../Code/OpenMV/FaceTracking.py}
	\caption{Gesichtsverfolgung \FILE{FaceTracking.py}} 	\label{MicroPythonFaceTracking}
\end{code}  



\subsection{Suche nach Kreisen um die Objekte}

In diesem Beispiel werden die Kreise in einem Bild mithilfe des Hough-Transform-Algorithmus gefunden. Dieser Algorithmus wird in der Bildverarbeitung zur Extraktion von Merkmalen wie Kreisen oder Ellipsen in einem Bild verwendet.  Um dies zu programmieren, wird der Sensor zunächst zurückgesetzt. Dann wird mit der Funktion \PYTHON{sensor.snapshot()} ein Schnappschuss vom Videostream gemacht. Dann wird der Kreis im Bild mit der Funktion \PYTHON{find\_circles()} gesucht. Zum Schluss werden Kreise mit der Funktion \PYTHON{draw\_circles()} gezeichnet.

\begin{figure}[H]
	\centering
	\includegraphics[width=\textwidth]{Arduino/find_circles}
	\caption{Ausgabe für das Finden von Kreisen}
	\label{OpenMVCircle}
\end{figure}

Wie in der Abbildung \ref{OpenMVCircle} zu sehen ist, zeigt die Ausgabe, dass die Kreise am Rand der Uhr gezeichnet werden. Dabei ist beachten, dass der Fokus der Kamera korrekt eingestellt wird. Da die Kamera eine feste Brennweite hat, muss der Abstand an die Objekte variiert werden. In der Regel müssen die Objekte näher an die Kamera gebracht werden, damit der Fokus richtig eingestellt ist, um die Kreise um die Objekte zu erkennen.
  

\begin{code}[H]
	\lstinputlisting[language=Python, firstnumber=1]{../Code/MicroPython/FindingCircle.py}
    \caption{Suche nach Kreisen um die Objekte \FILE{FindingCircle.py}} 	\label{MicroPythonCircle}
\end{code}  
 
  
\subsection{Extraktion von QR-Code-Informationen}

In diesem Beispiel die Informationen aus dem QR-Code mit Hilfe der Kamera extrahiert.  Dies ist eine interessante Funktion der OpenMV-Kamera, die im Vision Shield und im Necla Vision enthalten ist. Um dies zu programmieren, werden zunächst dieselben Initialisierungsschritten wie in den früheren Beispielen durchgeführt.  Dann wird ein Schnappschuss aus dem Videostream mit der Funktion \PYTHON{sensor.snapshot()} aufgenommen. Anschließend wird eine Objektivkorrektur auf 1,8 durchgeführt, um eine bessere Bildqualität mit der Funktion \PYTHON{lens\_corr()} zu erhalten.  Um den QR-Code zu extrahieren, steht die  Funktion \PYTHON{find\_qrcodes()} zur Verfügung. Abschließend wird ein Rechteck um den QR-Code mit der Funktion \PYTHON{draw\_rectangle()} gezeichnet.

\begin{figure}[H]
	\centering
	\includegraphics[width=\textwidth]{Arduino/QRcodedetection}
	\caption{Ausgabe der QR Code-Sketch}
	\label{OpenMVQR}
\end{figure}

In der Abbildung \ref{OpenMVQR}, die die Ausgabe der QR Code-Sketch darstellt, kann man sehen, dass ein rechteckiger Kasten um den QR-Code herum erstellt wird und die Informationen daraus extrahiert werden, die in der seriellen Ausgabestelle zu sehen sind. Das vollständige Programmcode ist in dem Programmbereich \ref{MicroPythonBarCode} dargestellt.

\begin{code}[H]
	\lstinputlisting[language=Python, firstnumber=1]{../Code/OpenMV/BarCode.py}
	\caption{Programm zur Extraktion von QR-Code-Informationen \FILE{BarCode.py}} 	\label{MicroPythonBarCode}
\end{code}  


